\scanpage{20}
\deleted{\uline{Corollaire 1 du th. 3} \inserted{Si $A$ est de dim $n$,} Le foncteur $M\unsure{-|\mapsto}H^n(M)$ est
exact à droite, et \deleted{\unsure{est}}\inserted{\unsure{°}} par suite}
\[
\deleted{\boxed{H^n(M)\xleftarrow{\sim}H^n(A)\otimes M}}
\]

\deleted{Corollaire 3 Si $D(\Omega)\simeq H^n(A)$, alors $H^n(\Omega)\simeq D(\operatorname{Hom}(\Omega,\Omega))$}

\inserted{\uline{Corollaire 2} \deleted{\unsure{Soit $M$}} \unsure{de} dim $q$, on a $\operatorname{Ext}^q(k,M)\ne0$. Si $i>q$ est tel que
$\operatorname{Ext}^i(k,M)\ne0$, alors $\operatorname{Ext}^j(k,M)=0$ pour $j>i$.}

\deleted{Soit \unsure{alors|donc} $\Omega$ donné, et $H^n(\Omega)\simeq\deleted{\unsure{D(\Omega)}}H^n(A)\otimes\Omega\xrightarrow{\varphi}\mathrm I$}

\deleted{\unsure{ce qui revient}}

\deleted{Il revient au \unsure{$=$|même} de se donner $\Omega\xrightarrow{\varphi_2}D(H^n(A))$
ou $H^n(A)\xrightarrow{\varphi_1}D(\Omega)$. On a alors des
accouplements}
\[
\deleted{H^p(M)\times\operatorname{Ext}^{n-p}(M,\Omega)\longrightarrow H^n(\Omega)\longrightarrow\mathrm I}
\]
\deleted{d'où}
\[
\deleted{H^p(M)\longrightarrow D\bigl(\operatorname{Ext}^{n-p}(M,\Omega)\bigr)}
\]

\deleted{Explicitons-les en dims extrêmes.}
\[
\deleted{H^n(M)\simeq\deleted{\unsure{M\otimes}}H^n(A)\otimes M\longrightarrow D(\operatorname{Hom}(M,\Omega))\simeq D(\Omega)\otimes M}
\]
\deleted{\deleted{est \unsure{précisément associé à} $\unsure{\varphi_1}$}
s'identifie à $\varphi_1\otimes\operatorname{id}_M$ ; pour que
ce soit un isom pour tt $M$, il f. et s.
que $\varphi_1$ soit un isom. Considérons le
cas \unsure{part.} \deleted{$\unsure{\varphi_1\otimes(M)}$} $M=\Omega$, on trouve}

% slanted
\margin{plus bas ?}

\[
\deleted{\deleted{H^n(\Omega)\unsure{\simeq} H^n(A)}}
\]
\[
\deleted{H^n(\Omega)\times\operatorname{Hom}(\Omega,\Omega)\longrightarrow H^n(\Omega)\longrightarrow\mathrm I}
\]
\deleted{d'où}
\[
\deleted{\widehat{\operatorname{Hom}(\Omega,\Omega)}\xrightarrow{\sim}D(H^n(\Omega))}
\]
\deleted{transformant l'identité en $\varphi$. \deleted{on} \uline{\unsure{donc|dira}}
\unsure{on a} équivalence des conditions suivantes :}

\deleted{(i) $H^n(\Omega)$ est dualisant et \unsure{que} \illegible{1}}

\deleted{(ii) $\operatorname{Hom}(\Omega,\Omega)\simeq A$ \qquad (iii)}
% archive sheet number: 19

